% !Mode:: "TeX:UTF-8"

%? ************************************************************* %?
%?                                                               %?
%?                   您将在此处完成致谢和附录内容                     %?
%?                                                               %?
%? ************************************************************* %?

% ====================== 致谢 ======================

\addcontentsline{toc}{chapter}{\texorpdfstring{\sihao\heiti{致\qquad 谢}}{致谢}}

\chapter*{\sanhao\heiti\textbf{致\qquad 谢}}

键盘敲下最后一个句号，屏幕上的文字终于定稿。回望这四年，有深夜调试代码的焦灼，也有仿真跑通那一瞬的畅快，此刻都化作了一段值得反复回味的记忆。

感谢我的指导老师刘鹏博老师。从最初选定仓储机器人这个方向，到剪叉机构的力学分析，再到Hybrid A*算法的反复调优，刘老师始终给予我方向上的指引和细节上的把关。每次带着一堆问题去请教，总能带着清晰的思路回来。刘老师身上那种对技术问题刨根问底的劲头，让我对"严谨"二字有了切身的理解。

感谢学长仪效哲。写论文的这几个月，学长在仿真调试和理论推导上给了我很多实质性的帮助，尤其是在PID参数整定和路径规划算法优化方面，少走了不少弯路。学长的耐心和慷慨，让我感受到了技术人之间最朴素的传承。

感谢我的几位舍友。四年的朝夕相处，从一起赶作业到期末互相提问，从深夜卧谈到互相打气，这些琐碎的日常构成了大学生活最温暖的底色。你们的陪伴让那些焦虑的夜晚不再难熬。

感谢我的父母。你们从不过问成绩排名，只关心我有没有好好吃饭、有没有照顾好自己。这种沉默而坚定的支持，是我在学业和生活中最坚实的后盾。希望这篇论文能让你们感到一丝欣慰。

所有的经历都是伏笔，所有的努力都不会被辜负。愿未来的路上，依然有代码可写，有难题可解，有热爱可循。


% ====================== 附录A：MATLAB仿真程序代码 ======================

\chapter*{\sanhao\heiti\textbf{附录A\quad MATLAB仿真程序代码}}
\addcontentsline{toc}{chapter}{\texorpdfstring{\sihao\heiti{附录A\quad MATLAB仿真程序代码}}{附录A}}

\noindent \textbf{附A1}\quad 典型轨迹运动学仿真程序 \texttt{traj\_sim.m}

\lstinputlisting[language=Matlab]{static/code/traj_sim.m}

\vspace{1em}

\noindent \textbf{附A2}\quad Hybrid A*路径规划仿真程序 \texttt{hybrid\_astar.m}

\lstinputlisting[language=Matlab]{static/code/hybrid_astar.m}

\vspace{1em}

\noindent \textbf{附A3}\quad 红外循迹PID仿真程序 \texttt{tracking\_pid.m}

\lstinputlisting[language=Matlab]{static/code/tracking_pid.m}

\vspace{1em}

\noindent \textbf{附A4}\quad 双环PID运动控制仿真程序 \texttt{pid\_control.m}

\lstinputlisting[language=Matlab]{static/code/pid_control.m}

\vspace{1em}

\noindent \textbf{附A5}\quad 剪叉式升降机构强度校核与载重分析程序 \texttt{lift\_strength.m}

\lstinputlisting[language=Matlab]{static/code/lift_strength.m}



